The relationship between internal neural computations and non-invasive recordings like electroencephalography (EEG) remains poorly understood. EEG provides a window into large-scale brain activity, yet the mapping from cellular/circuit-level computations to surface recordings is complex, involving spatial mixing, filtering, and volume conduction effects.

A central question in neuroscience is: \textit{What can we infer about internal computational structures from non-invasive recordings?} This question is particularly acute given the growing use of deep neural networks as models of neural computation \cite{Kriegeskorte2015, Yamins2016}. If artificial neural networks capture principles of biological computation, understanding how internal routing structures project to observable signals could illuminate the neural basis of cognition.

Recent work has highlighted the importance of sparse attention in Transformers, where selective focus on relevant information improves both computational efficiency and interpretability \cite{Child2019, Zaheer2020}. Yet it remains unclear whether changes in internal attention routing produce measurable signals at the population level.

We approach this as a causal inference problem: \textit{If we experimentally control attention sparsity in a neural network and observe corresponding changes in downstream signals, how well does the inferred causal structure transfer to real neural recordings?}

This paper tests this through a systematic protocol:

\begin{enumerate}
    \item Create a controlled system (transformer with sparse attention)
    \item Extract pathway metrics (6 routing measures from attention)
    \item Project to EEG-like format (linear mapping to 105 channels)
    \item Quantify information retention through predictive modeling
    \item Test causal transfer (pathway $\rightarrow$ EEG generalization)
\end{enumerate}

The goal is not to model real EEG, but to test whether a \textit{principle}—"computational routing structure projects to population signals"—holds under controlled conditions. Success would suggest that EEG interpretability through the lens of internal routing is theoretically grounded.
